\documentclass[10pt, letterpaper]{article}

% Packages:
\usepackage[
    ignoreheadfoot, % set margins without considering header and footer
    top=2 cm, % seperation between body and page edge from the top
    bottom=2 cm, % seperation between body and page edge from the bottom
    left=2 cm, % seperation between body and page edge from the left
    right=2 cm, % seperation between body and page edge from the right
    footskip=1.0 cm, % seperation between body and footer
    % showframe % for debugging 
]{geometry} % for adjusting page geometry
\usepackage{titlesec} % for customizing section titles
\usepackage{tabularx} % for making tables with fixed width columns
\usepackage{array} % tabularx requires this
\usepackage[dvipsnames]{xcolor} % for coloring text
\definecolor{primaryColor}{RGB}{0, 0, 0} % define primary color
\usepackage{enumitem} % for customizing lists
\usepackage{fontawesome5} % for using icons
\usepackage{amsmath} % for math
\usepackage[
    pdftitle={Guangzhi Xiong's CV},
    pdfauthor={Guangzhi Xiong},
    pdfcreator={LaTeX with RenderCV},
    colorlinks=true,
    urlcolor=primaryColor
]{hyperref} % for links, metadata and bookmarks
\usepackage[pscoord]{eso-pic} % for floating text on the page
\usepackage{calc} % for calculating lengths
\usepackage{bookmark} % for bookmarks
\usepackage{lastpage} % for getting the total number of pages
\usepackage{changepage} % for one column entries (adjustwidth environment)
\usepackage{paracol} % for two and three column entries
\usepackage{ifthen} % for conditional statements
\usepackage{needspace} % for avoiding page brake right after the section title
\usepackage{iftex} % check if engine is pdflatex, xetex or luatex

% Ensure that generate pdf is machine readable/ATS parsable:
\ifPDFTeX
    \input{glyphtounicode}
    \pdfgentounicode=1
    \usepackage[T1]{fontenc}
    \usepackage[utf8]{inputenc}
    \usepackage{lmodern}
\fi

\usepackage{charter}

\usepackage{graphicx}
\usepackage{hyperref}

% Some settings:
\raggedright
\AtBeginEnvironment{adjustwidth}{\partopsep0pt} % remove space before adjustwidth environment
\pagestyle{empty} % no header or footer
\setcounter{secnumdepth}{0} % no section numbering
\setlength{\parindent}{0pt} % no indentation
\setlength{\topskip}{0pt} % no top skip
\setlength{\columnsep}{0.15cm} % set column seperation
\pagenumbering{gobble} % no page numbering

\titleformat{\section}{\needspace{4\baselineskip}\bfseries\large}{}{0pt}{}[\vspace{1pt}\titlerule]

\titlespacing{\section}{
    % left space:
    -1pt
}{
    % top space:
    0.3 cm
}{
    % bottom space:
    0.2 cm
} % section title spacing

\renewcommand\labelitemi{$\vcenter{\hbox{\small$\bullet$}}$} % custom bullet points
\newenvironment{highlights}{
    \begin{itemize}[
        topsep=0.10 cm,
        parsep=0.10 cm,
        partopsep=0pt,
        itemsep=0pt,
        leftmargin=0 cm + 10pt
    ]
}{
    \end{itemize}
} % new environment for highlights


\newenvironment{highlightsforbulletentries}{
    \begin{itemize}[
        topsep=0.10 cm,
        parsep=0.10 cm,
        partopsep=0pt,
        itemsep=0pt,
        leftmargin=10pt
    ]
}{
    \end{itemize}
} % new environment for highlights for bullet entries

\newenvironment{onecolentry}{
    \begin{adjustwidth}{
        0 cm + 0.00001 cm
    }{
        0 cm + 0.00001 cm
    }
}{
    \end{adjustwidth}
} % new environment for one column entries

\newenvironment{twocolentry}[2][]{
    \onecolentry
    \def\secondColumn{#2}
    \setcolumnwidth{\fill, 4.5 cm}
    \begin{paracol}{2}
}{
    \switchcolumn \raggedleft \secondColumn
    \end{paracol}
    \endonecolentry
} % new environment for two column entries

\newenvironment{threecolentry}[3][]{
    \onecolentry
    \def\thirdColumn{#3}
    \setcolumnwidth{, \fill, 4.5 cm}
    \begin{paracol}{3}
    {\raggedright #2} \switchcolumn
}{
    \switchcolumn \raggedleft \thirdColumn
    \end{paracol}
    \endonecolentry
} % new environment for three column entries

\newenvironment{header}{
    \setlength{\topsep}{0pt}\par\kern\topsep\centering\linespread{1.5}
}{
    \par\kern\topsep
} % new environment for the header

\newcommand{\placelastupdatedtext}{% \placetextbox{<horizontal pos>}{<vertical pos>}{<stuff>}
  \AddToShipoutPictureFG*{% Add <stuff> to current page foreground
    \put(
        \LenToUnit{\paperwidth-2 cm-0 cm+0.05cm},
        \LenToUnit{\paperheight-1.0 cm}
    ){\vtop{{\null}\makebox[0pt][c]{
        \small\color{gray}\textit{Last updated in September 2024}\hspace{\widthof{Last updated in September 2024}}
    }}}%
  }%
}%

% save the original href command in a new command:
\let\hrefWithoutArrow\href

% new command for external links:

% \immediate\write18{curl -L -o medrag-stars.png "https://img.shields.io/github/stars/Teddy-XiongGZ/MedRAG?style=social"}


\begin{document}
    \newcommand{\AND}{\unskip
        \cleaders\copy\ANDbox\hskip\wd\ANDbox
        \ignorespaces
    }
    \newsavebox\ANDbox
    \sbox\ANDbox{$|$}

    \begin{header}
        \fontsize{25 pt}{25 pt}\selectfont Lei Gong

        \vspace{5 pt}

        \normalsize
        \mbox{Address: Rice Hall 542, 85 Engineer's Way, Charlottesville, Virginia 22904}%
        \kern 5.0 pt%
        \AND%
        \kern 500.0 pt%
        \mbox{\hrefWithoutArrow{nkw3mr@virginia.edu}{nkw3mr@virginia.edu}}%
        \kern 5.0 pt%
        \AND%
        \kern 5.0 pt%
        \mbox{\hrefWithoutArrow{tel:+1-765-771-9032}{(+1) 765-771-9032}}%
        % \kern 5.0 pt%
        % \AND%
        % \kern 5.0 pt%
        % \mbox{\hrefWithoutArrow{https://yourwebsite.com/}{yourwebsite.com}}%
        \kern 5.0 pt%
        \AND%
        \kern 5.0 pt%
        \mbox{\hrefWithoutArrow{https://github.com/LeiGong0125Carrot?tab=repositories}{GitHub}}%
        \kern 5.0 pt%
        \AND%
        \kern 5.0 pt%
        \mbox{\hrefWithoutArrow{https://scholar.google.com/citations?user=8QLAkKMAAAAJ&hl=zh-CN}{Google Scholar}}%

    \end{header}

    \vspace{5 pt - 0.3 cm}


    % \section{Welcome to RenderCV!}
    %     \begin{onecolentry}
    %         \href{https://rendercv.com}{RenderCV} is a LaTeX-based CV/resume version-control and maintenance app. It allows you to create a high-quality CV or resume as a PDF file from a YAML file, with \textbf{Markdown syntax support} and \textbf{complete control over the LaTeX code}.
    %     \end{onecolentry}

    %     \vspace{0.2 cm}

    %     \begin{onecolentry}
    %         The boilerplate content was inspired by \href{https://github.com/dnl-blkv/mcdowell-cv}{Gayle McDowell}.
    %     \end{onecolentry}


    % \section{Quick Guide}
    % \begin{onecolentry}
    %     \begin{highlightsforbulletentries}
    %     \item Each section title is arbitrary and each section contains a list of entries.
    %     \item There are 7 unique entry types: \textit{BulletEntry}, \textit{TextEntry}, \textit{EducationEntry}, \textit{ExperienceEntry}, \textit{NormalEntry}, \textit{PublicationEntry}, and \textit{OneLineEntry}.
    %     \item Select a section title, pick an entry type, and start writing your section!
    %     \item \href{https://docs.rendercv.com/user_guide/}{Here}, you can find a comprehensive user guide for RenderCV.
    %     \end{highlightsforbulletentries}
    % \end{onecolentry}
    \section{RESEARCH INTERESTS}
My research focuses on developing reliable and knowledge-enhanced machine learning methods for healthcare, with an emphasis on patient-level outcome prediction, clinical representation learning, and learning from heterogeneous and incomplete clinical data. I am particularly interested in bridging clinical NLP, foundation models, and multimodal health intelligence to support robust and individualized healthcare decision-making.
    
   \section{EDUCATION}

\begin{twocolentry}{
    United States
}
    \textbf{University of Virginia}
\end{twocolentry}

\vspace{0.05 cm}
\begin{onecolentry}
    \begin{highlights}
        \item \begin{twocolentry}{
            01/2024 -- Present
        }
            Master in \textit{Computer Science}
        \end{twocolentry}
        \begin{highlights}
            \item GPA: 3.91/4.00
            \item Advisor: \href{https://www.cs.virginia.edu/~az9eg/website/home.html}{Aidong Zhang}
        \end{highlights}
    \end{highlights}
\end{onecolentry}

\vspace{0.10 cm}

\begin{twocolentry}{
    United States
}
    \textbf{Purdue University}
\end{twocolentry}

\vspace{0.05 cm}
\begin{onecolentry}
    \begin{highlights}
        \item \begin{twocolentry}{
            01/2019 -- 05/2023
        }
            \textbf{Bachelor of Science} in \textit{Computer Science}; 
            \textbf{Second Major} in \textit{Applied Statistics}
        \end{twocolentry}
        \begin{highlights}
            \item GPA: 3.55/4.00
            \item \textbf{Honors/Awards:} Dean's List and Semester Honors, Fall 2021 and Spring 2022; Semester Honor, Spring 2020.
        \end{highlights}
    \end{highlights}
\end{onecolentry}
            
    

        
    \section{PUBLICATIONS}
        \begin{enumerate}[label={[\arabic*]}]
            
            \item \textbf{Boosting Clinical Outcome Prediction with Context-Aware Feature Imputation and Disentanglement} [\href{https://ieeexplore.ieee.org/abstract/document/11392038}{paper}] [\href{https://github.com/LeiGong0125Carrot/Strucure-Awared-Clinical-Note-Processing}{code}] \\
            \underline{Lei Gong}, Aidong Zhang, Kishlay Jha \\
            \textbf{ICDM}, 2025
            
            \item \textbf{Context-specific Feature Augmentation for Improving Social Determinants of Health Extraction} [\href{https://ieeexplore.ieee.org/abstract/document/10825225}{paper}] \\
            \underline{Lei Gong}, Andrey Shor, Aidong Zhang, Kishlay Jha  \\
            \textbf{Big Data}, 2024
            
            \item \textbf{Boosting Social Determinants of Health Extraction with Semantic Knowledge Augmented Large Language Model} [\href{https://pmc.ncbi.nlm.nih.gov/articles/PMC12099417/}{paper}]\\
            \underline{Lei Gong}, Jaren Bresnick, Aidong Zhang, Cathy Wu, Kishlay Jha  \\
            \textbf{AMIA}, 2024

            \item \textbf{SCoPe: Semantic Context-Adaptive Prototypical Learning for Clinical
Outcome Predictio} [\href{https://github.com/LeiGong0125Carrot/clinical_diffusion_generator/tree/concept_merging_vib}{code}] \\
            \underline{Lei Gong}, Kanghui Ning, Kishlay Jha, Aidong Zhang \\
            

            \item \textbf{Reward Is Enough: LLMs Are In-Context Reinforcement Learners} [\href{https://arxiv.org/abs/2506.06303}{paper}] [\href{https://github.com/garified/ICRL-for-LLM-Agent}{code}] \\
            Kefan Song, Amir Moeini, Peng Wang, \underline{Lei Gong}, Rohan Chandra, Shangtong Zhang, Yanjun Qi\\
            \textbf{ICLR}, 2026

            \item \textbf{OPD-Aha: From Linguistic Momentum to Visual Reflection in Multimodal On-Policy Distillation} [\href{https://arxiv.org/abs/2609.16459}{paper}] [\href{https://github.com/Echochef/OPD-Aha}{code}] \\
            Chenhao Qiu, Dawei Li, Yechao Zhang, \underline{Lei Gong}, Zhen Tan\\
            \textbf{arXiv}, 2026
                    
        \end{enumerate}

        \section{RESEARCH EXPERIENCE}
        \textbf{Graduate Research Assistant, University of Virginia} \hfill Charlottesville, VA \\
Advisor: Prof. Aidong Zhang \hfill 01/2024 -- Present
\begin{itemize}
    \item Conducted research on \textbf{interpretable clinical NLP} and reliable EHR-based prediction, developing \textbf{knowledge- and retrieval-augmented learning} methods for \textbf{social determinants of health extraction}, \textbf{context-aware feature augmentation and imputation} methods for semi-structured clinical narratives, and \textbf{clinical knowledge graphs} from large-scale clinical notes.
\end{itemize}
        
\section{SELECT PROJECTS}


% 1. SCoPe
\begin{twocolentry}{
    2026 -- Present
}
    \textbf{SCoPe: Semantic Context-Adaptive Prototypical Learning for Clinical Outcome Prediction}
\end{twocolentry}

\vspace{0.05 cm}
\begin{onecolentry}
    \begin{highlights}
        \item Developed a \textbf{context-adaptive representation learning} framework with dynamic semantic unit aggregation, class-conditional prototypical learning, and KL regularization for interpretable clinical outcome prediction.

        \item Outperformed strong baselines on MIMIC-III, MIMIC-IV, and Hyperpartisan benchmarks, improving mortality prediction F1 by \textbf{14.9\%} (0.268 to 0.308) over Clinical-ModernBERT.
    \end{highlights}
\end{onecolentry}

\vspace{0.10 cm}

% 2. OPD-Aha
\begin{twocolentry}{
    2026 -- Present
}
    \textbf{OPD-Aha: From Linguistic Momentum to Visual Reflection in Multimodal On-Policy Distillation}
\end{twocolentry}

\vspace{0.05 cm}
\begin{onecolentry}
    \begin{highlights}
        \item Collaborated on \textbf{multimodal on-policy distillation} research, contributing to algorithm development discussions and \textbf{Qwen3.5-4B/9B post-training} with visual preference-based target reconstruction.

        \item Contributed to experimental validation and benchmark evaluation, with OPD-Aha achieving \textbf{+4.0/+3.2 percentage-point} average accuracy gains over Vision-OPD across six visual reasoning benchmarks.
    \end{highlights}
\end{onecolentry}

\vspace{0.10 cm}

% 3. ICDM 2025

\begin{twocolentry}{
    ICDM 2025
}
    \textbf{Boosting Clinical Outcome Prediction with Context-Aware Feature Imputation and Disentanglement}
\end{twocolentry}

\vspace{0.05 cm}
\begin{onecolentry}
    \begin{highlights}
        \item Developed a \textbf{context-aware feature imputation} framework combining structure-guided multi-head attention and \textbf{disentangled representation learning} to improve clinical outcome prediction from incomplete EHR data.

        \item Evaluated on \textbf{MIMIC-III} clinical benchmarks, improving macro-F1 by \textbf{4.44 percentage points} for prolonged mechanical ventilation prediction (COReBERT) and \textbf{1.81 percentage points} for mortality prediction (BlueBERT).
    \end{highlights}
\end{onecolentry}


\vspace{0.10 cm}

% 4. Big Data 2024

\begin{twocolentry}{
    Big Data 2024
}
    \textbf{Context-Specific Feature Augmentation for Social Determinants of Health Extraction}
\end{twocolentry}

\vspace{0.05 cm}
\begin{onecolentry}
    \begin{highlights}
        \item Developed a \textbf{retrieval-augmented clinical NLP} framework integrating biomedical literature with EHR narratives through \textbf{cross-source attention and gated feature fusion} for knowledge-enhanced representation learning.

        \item Outperformed six baselines on the \textbf{MIMIC-IV SDoH} benchmark, improving \textbf{macro-F1 by 2.91 percentage points} (8.53\% to 11.44\%) and \textbf{macro-AUC by 2.49 percentage points} (72.92\% to 75.41\%) over Hierarchical RoBERTa.
    \end{highlights}
\end{onecolentry}

\vspace{0.10 cm}


% 5. AMIA 2024
\begin{twocolentry}{
    AMIA 2024
}
    \textbf{Boosting Social Determinants of Health Extraction with Semantic Knowledge Augmented Large Language Model}
\end{twocolentry}

\vspace{0.05 cm}
\begin{onecolentry}
    \begin{highlights}
        \item Developed a \textbf{knowledge-augmented language model} framework integrating UMLS biomedical semantic types into clinical narratives to address data scarcity and class imbalance in SDoH extraction.

        \item Evaluated \textbf{FLAN-T5 models} across multiple scales, improving macro-F1 by \textbf{4.01 percentage points} (0.8423 to 0.8824) on MIMIC-III SDoH data.
    \end{highlights}
\end{onecolentry}

\vspace{0.10 cm}















    \section{TEACHING \& SERVICES}
    % \section{Teaching}
    
    \begin{onecolentry}
    \textbf{Teaching Assistant} 
    \begin{highlights}
    \item Purdue POL 300 - Introduction to Political Analysis (Fall 2022 - Spring 2023)
    \end{highlights}
    \end{onecolentry}

    % \section{Services}
    \vspace{0.2 cm}
    
    \begin{onecolentry}
    \textbf{Conference Reviewer}
    \begin{highlights}
    \item SDM 2024, ICDM 2025
    \end{highlights}
    \end{onecolentry}
    
    % \section{Technical Skills}

    %     \begin{onecolentry}
    %         \textbf{Programming:} Python, R, C/C++
    %     \end{onecolentry}

    %     \vspace{0.2 cm}

    %     \begin{onecolentry}
    %         \textbf{Tools \& Platforms:} PyTorch, Transformers, LangChain, MedRAG
    %     \end{onecolentry}
    \section{TECHNICAL SKILLS}

\begin{onecolentry}
\textbf{Languages:} Python, SQL, R, Java, C \\
\textbf{ML/NLP:} Pytorch, Clinical NLP, information extraction, multi-label prediction, large language models, transformer-based modeling, MetaMap \\
\textbf{Data/Knowledge Systems:} Neo4j, knowledge graph construction \\
\textbf{Visualization:} Tableau \\
\textbf{Languages:} English (Fluent), Mandarin (Native)
\end{onecolentry}

    
    

    

\end{document}